\documentclass[10pt,a4paper]{amsart}
\usepackage[margin=19mm]{geometry}
\usepackage{amsmath,amssymb,mathtools}
\usepackage[scr=rsfs,cal=euler]{mathalfa}
\usepackage{xfrac}
\usepackage[unicode,hidelinks,bookmarks=false]{hyperref}
\newcommand{\ie}{i.e.\ }
\newcommand{\cf}{cf.\ }
\newcommand{\resp}{respectively\ }
\newcommand{\Subsection}{\subsection}
\providecommand{\nobreakdash}{\nobreak\hbox{-}}
\setlength{\emergencystretch}{2em}
\multlinegap0pt
\usepackage{bm}
\usepackage{bbm}
\usepackage{dsfont}
\usepackage{xspace}
\usepackage{cancel}
\usepackage{mathtools}
\usepackage{amsthm}
\makeatletter
\@ifundefined{theorem}{\newtheorem{theorem}{Theorem}}{}
\@ifundefined{definition}{\newtheorem{definition}[theorem]{Definition}}{}
\@ifundefined{proposition}{\newtheorem{proposition}[theorem]{Proposition}}{}
\makeatother
\DeclareMathOperator{\cw}{cw}
\newcommand\ba{\mathbf{a}}
\newcommand\bb{\mathbf{b}}
\newcommand{\cliques}{\mathcal{C}}
\newcommand{\cliquevs}{\mathcal{V}}
\title[Flow polytopes for extensions of bipartite graphs]{Flow polytopes for extensions of bipartite graphs}
\author{Benjamin Braun, Kaitlin Bruegge, Robert Davis, Derek Hanely}
\date{Source: arXiv:2509.26445v2 (CC BY 4.0)}
\begin{document}
\maketitle

\noindent\emph{Source and licence.} Flow polytopes for extensions of bipartite graphs. Version arXiv:2509.26445v2 (CC BY 4.0), distributed under the Creative Commons Attribution 4.0 International licence, as recorded in the arXiv licence metadata for this exact version and shown on the abstract page https://arxiv.org/abs/2509.26445v2. Full rights notice: licenses/arxiv-2509.26445-CC-BY-4.0.txt.\par\smallskip

\noindent\textit{Editorial scope.} Counted unit: the Proposition whose source label is \texttt{prop:cover} in the article (source0.tex lines 1203--1226, source label \texttt{prop:cover}), delivered together with the article's own complete original proof of it (source0.tex lines 1228--1311, one \texttt{proof} environment of 84 lines). No mathematics is written, repaired or re-derived: statement and proof are the article's own text line for line. Required context retained uncounted: def:cliquemap (lines 730--754). Every definition, display and label of those lines is retained unchanged. Omitted from this package and not redistributed: 1--729, 755--1202, 1227--1227, 1312--1447. Nothing inside a retained excerpt is omitted or altered: every retained body is delivered exactly as the article's lines are written, so no omission receipt of any kind accompanies this package. Citations: the retained excerpts carry no citation command at all, so no bibliography entry of the article is transcribed and no reference is redistributed. Reference adaptation: none was needed and none was made; every reference inside the delivered unit resolves inside it, so no unresolved reference or citation is produced. Printed slips kept literal: no printed word, symbol, hypothesis or number of the counted statement or of its proof was altered. Layout and class: the article's own class is \texttt{amsart} with packages including amssymb, mathrsfs, xy, enumitem, natbib, url, verbatim, comment, geometry, mathtools, caption, rotating, tikz, graphicx; the wrapper is the lane's pinned \texttt{amsart} editorial wrapper at 10pt on A4, which restates the theorem-like environments the retained text needs and adds the article's own macro definitions that the retained text needs (\texttt{\textbackslash ba}, \texttt{\textbackslash bb}, \texttt{\textbackslash cliques}, \texttt{\textbackslash cliquevs}, \texttt{\textbackslash cw}), with extra wrapper packages (bm, bbm, dsfont, xspace, cancel, mathtools). The delivered numbering is the unit's own. Assets: no retained span contains a figure environment, a graphics-inclusion command, a PDF-page inclusion, a table or a TikZ picture; no image file is copied into this package, the package declares no asset dependency and no image file is redistributed with it. Licence: version arXiv:2509.26445v2 of this article is distributed under the Creative Commons Attribution 4.0 International licence, as recorded in the arXiv licence metadata for this exact version and shown on the abstract page https://arxiv.org/abs/2509.26445v2. A full rights notice is delivered at licenses/arxiv-2509.26445-CC-BY-4.0.txt. Characters: every retained excerpt is pure ASCII, so the delivered engine, XeLaTeX with its default Latin Modern text font, reports no missing character.
\begin{definition}\label{def:cliquemap}
    We define the \emph{clique map} $\phi:\cliquevs(H)\to \cliques(H)$ as follows.
    For each edge $ij\in E(H)$, we include in $\phi(\ba)$ the following routes:
    \begin{enumerate}
    \item[(i)] $a_i\neq j$ and $a_j\neq i$:
        \begin{itemize}
        \item If $a_i<j$ and $a_j<i$, include the route $\alpha_{2,i}\beta_{i,j}\gamma_{j,2}$.
        \item If $a_i<j$ and $a_j>i$, include the route $\alpha_{2,i}\beta_{i,j}\gamma_{j,1}$.
        \item If $a_i>j$ and $a_j<i$, include the route $\alpha_{1,i}\beta_{i,j}\gamma_{j,2}$.
        \item If $a_i>j$ and $a_j>i$, include the route $\alpha_{1,i}\beta_{i,j}\gamma_{j,1}$.
        \end{itemize}
        \item[(ii)] Either $a_i=j$ or $a_j=i$, but not both:
        \begin{itemize}
        \item If $a_i=j$ and $a_j<i$, include the routes $\alpha_{1,i}\beta_{i,j}\gamma_{j,2}$ and $\alpha_{2,i}\beta_{i,j}\gamma_{j,2}$.
        \item If $a_i=j$ and $a_j>i$, include the routes $\alpha_{1,i}\beta_{i,j}\gamma_{j,1}$ and $\alpha_{2,i}\beta_{i,j}\gamma_{j,1}$.
        \item If $a_j=i$ and $a_i<j$, include the routes $\alpha_{2,i}\beta_{i,j}\gamma_{j,1}$ and $\alpha_{2,i}\beta_{i,j}\gamma_{j,2}$.
        \item If $a_j=i$ and $a_i>j$, include the routes $\alpha_{1,i}\beta_{i,j}\gamma_{j,1}$ and $\alpha_{1,i}\beta_{i,j}\gamma_{j,2}$.
        \end{itemize}
        \item[(iii)] $a_i=j$ and $a_j=i$:
    \begin{itemize}
        \item If $a_{i,j}=-$, include the routes $\alpha_{1,i}\beta_{i,j}\gamma_{j,1}$, $\alpha_{1,i}\beta_{i,j}\gamma_{j,2}$, and $\alpha_{2,i}\beta_{i,j}\gamma_{j,2}$.
        \item If $a_{i,j}=+$, include the routes $\alpha_{1,i}\beta_{i,j}\gamma_{j,1}$, $\alpha_{2,i}\beta_{i,j}\gamma_{j,1}$, and $\alpha_{2,i}\beta_{i,j}\gamma_{j,2}$.
        \end{itemize}
    \end{enumerate}
\end{definition}

\begin{proposition}\label{prop:cover}
    Let $\ba, \bb$ be clique vectors of $G(H)$ with respect to the canonical bipartite framing.
    Then $\phi(\ba) \lessdot \phi(\bb)$ if and only if exactly one of the following hold:
    \begin{enumerate}
        \item $\ba$ and $\bb$ differ in exactly one entry $a_i \neq b_i$, and
            \begin{enumerate}
                \item $i \in \{1,\dots,n\}$,
                \item $a_i = j$ and $a_j \neq i$, and
                \item $b_i$ is the largest neighbor of $i$ less than $j$.
            \end{enumerate}
        \item $\ba$ and $\bb$ differ in exactly one entry $a_j \neq b_j$, and
            \begin{enumerate}
                \item $j \in \{n+1,\dots,n+m\}$,
                \item $a_i \neq j$ and $a_j = i$, and
                \item $b_j$ is the smallest neighbor of $j$ greater than $i$.
            \end{enumerate}
        \item $\ba$ and $\bb$ differ only in that $a_{i,j} = -$ while $b_{i,j} = +$.
        \item $\ba$ and $\bb$ differ in exactly two entries as follows:
            \begin{enumerate}
                \item $a_{i,j} = +$ while $b_{i,j} = -$, and
                \item either $b_i$ is the largest neighbor of $i$ that is less than $j$ or $b_j$ is the smallest neighbor of $j$ that is greater than $i$.
            \end{enumerate}
    \end{enumerate}
\end{proposition}

\begin{proof}
    First suppose $\phi(\ba) \lessdot \phi(\bb)$ and let $S = \phi(\ba) \cap \phi(\bb)$.
    This means $\phi(\ba)$ and $\phi(\bb)$ are maximal cliques of $G(H)$ such that
    \begin{equation}\label{eq: clique difference}
       S = \phi(\ba) \setminus \{R_\ba\} = \phi(\bb) \setminus \{R_\bb\}
    \end{equation}
    for some routes $R_\ba \in \phi(\ba)$ and $R_\bb \in \phi(\bb)$.
    These routes satisfy $R_\ba <^{\cw} R_\bb$.
    This can occur in three ways:
    \begin{enumerate}
        \item[(A)] $R_\ba = \alpha_{1,i}\beta_{i,j}\gamma_{j,2}$ and $R_\bb = \alpha_{2,i}\beta_{i,j}\gamma_{j,1}$ for some edge $ij \in H$;
        \item[(B)] $R_\ba = \alpha_{1,i}\beta_{i,j'}\gamma_{j',*}$ and $R_\bb = \alpha_{2,i}\beta_{i,j}\gamma_{j,*}$ for some $n+1 \leq j < j' \leq n+m$; or
        \item[(C)] $R_\ba = \alpha_{*,i}\beta_{i,j}\gamma_{j,2}$ and $R_\bb = \alpha_{*,i'}\beta_{i',j}\gamma_{j,1}$ for some $1 \leq i < i' \leq n$. 
    \end{enumerate}

    \textbf{Case (A)}: if $\beta_{i,j}$ is used once in $\phi(\ba)$ and $\phi(\bb)$, then $a_i > j$ and $b_i < j$.
    However, by \eqref{eq: clique difference}, this means the set of routes passing through $\beta_{i,a_i}$ is the same in $\phi(\ba)$ and $\phi(\bb)$, hence $a_i = b_i$, which is impossible. 
    
    Now suppose $\beta_{i,j}$ is used twice in each of $\phi(\ba)$ and $\phi(\bb)$.
    By \eqref{eq: clique difference}, either both $a_i = j$ and $a_j< i$ or both $a_i>j$ and and $a_j = i$. 
    Without loss of generality, suppose the latter. 
    This means $b_i = j$ and $b_j > i$.
    However, this means $\beta_{i,a_i}$ is used a different number of times in $\phi(\ba)$ and in $\phi(\bb)$, contradicting \eqref{eq: clique difference}.

    Lastly, suppose $\beta_{i,j}$ appears three times in each of $\phi(\ba)$ and $\phi(\bb)$.
    Then $a_i = b_i = j$, $a_j = b_j = i$, and $a_{i,j} = -b_{i,j} = -$.
    By definition, both $\ba$ and $\bb$ are valid clique vectors.
    This falls into condition \emph{(3)} of the proposition statement.

    \textbf{Case (B)}: Since $\phi(\bb)$ is obtained from $\phi(\ba)$, $\beta_{i,j}$ is part of at most two routes in $\phi(\ba)$ and $\beta_{i,j'}$ is part of at least two routes in $\phi(\ba)$.
    First suppose $\beta_{i,j}$ appears in one route in $\phi(\ba)$ and $\beta_{i,j'}$ is in two.
    One way this can happen is if $a_i > j'$ and $a_{j'} = i$.
    When this happens, $b_{j'} \neq i$.
    However, \eqref{eq: clique difference} implies $b_{j'} \neq k$ for all $k \neq i$ as well, so this cannot be the case.

    The other way in which $\beta_{i,j}$ appears once and $\beta_{i,j'}$ appears twice is when $a_i = j'$ and $a_{j'} \neq i$.
    In each of the four cases of $a_j,a_{j'}$ being greater than or less than $i$, we find that $b_i = j$, $b_k = a_k$ for all $k \neq i$, and $b_{k,l} = a_{k,l}$ for all $k,l$.
    As a result, $\bb$ is a valid clique vector.

    In order for $\phi(\bb)$ to cover $\phi(\ba)$, we must have that $j$ is the largest neighbor of $i$ that is less than $j'$. 
    This falls into condition \emph{(1)} of the statement of the proposition.

    Next suppose that $\beta_{i,j}$ and $\beta_{i,j'}$ each appear in exactly two routes in $\phi(\ba)$. 
    This occurs when $a_i = b_i = j$, $b_j = i$, $a_j \neq i$, and $a_{j'} = i$.
    As a result, one route in $\phi(\ba)$ containing $\beta_{i,j}$ must begin with $\alpha_{1,j}$ and the other begins with $\alpha_{2,i}$, and both routes containing $\beta_{i,j'}$ must both begin with $\alpha_{2,i}$.
    However, this does not allow any route to be chosen as $R_\ba$ in a way that is consistent with the conditions assumed in case (B).
    So, this cannot occur.

    Now suppose $\beta_{i,j}$ is used in one route of $\phi(\ba)$ and $\beta_{i,j'}$ is used in three.
    This means $a_i = j'$, $a_{j'} = i$, and $a_j \neq i$.
    In order for $R_\bb$ to be coherent with the remaining routes passing through $i$, it must hold that $a_{i,j'} = +$, $b_i < j'$, and $b_{j'} = i$.
    Consequently, $b_i = j$ and $b_j = a_j$, forcing $b_{i,j'} = -$ and all other entries of $\bb$ are the same as those of $\ba$. 
    For $\phi(\bb)$ to cover $\phi(\ba)$, $j$ must be the largest neighbor of $i$ that is smaller than $j'$.
    It follows that $\bb$ is a valid clique vector, falling into condition \emph{(4)} of the proposition statement.

    Lastly, if $\beta_{i,j}$ is used twice and $\beta_{i,j'}$ is used three times in $\phi(\ba)$, then $a_i = j'$, $a_j = a_{j'} = i$, and $a_{i,j'} = +$.
    This forces $b_i = j$, $b_j = b_{j'} = i$, and $b_{i,j} = -$.
    This again results in a $\bb$ being a valid clique vector and falling into condition \emph{(4)} of the proposition statement. 

    \textbf{Case (C)}: This case is symmetric to Case (B).
    In particular, this case accounts for condition \emph{(2)} in the statement of the proposition and for the remaining part of condition \emph{(4)}.

    To establish the converse, suppose $\ba$ and $\bb$ satisfy one of the conditions in the statement of the proposition.
    The argument for verifying $\phi(\ba) \lessdot \phi(\bb)$ is similar enough for each condition, so we only present the argument for one of them.
    
    Suppose condition \emph{(1)} is satisfied, and without loss of generality suppose $a_j < i$.
    By Definition~\ref{def:cliquemap}, the clique $\phi(\ba)$ will contain the routes
    \[
        \alpha_{1,i}\beta_{i,j}\gamma_{j,2}~\text{ and }~\alpha_{2,i}\beta_{i,j}\gamma_{j,2},
    \]
    and $\phi(\bb)$ will only contain $\alpha_{2,i}\beta_{i,j}\gamma_{j,2}$.
    It is possible to have $a_k = a_{b_i} < i$, $a_k = i$, or $a_k > i$; again without loss of generality, suppose $a_k < i$.
    This means the only route containing $\beta_{i,k}$ in $\phi(\ba)$ is $\alpha_{1,i}\beta_{i,k}\gamma_{k,2}$, while $\phi(\bb)$ contains the routes
    \[
        \alpha_{1,i}\beta_{i,k}\gamma_{k,2}~\text{ and }~\alpha_{2,i}\beta_{i,k}\gamma_{k,2}.
    \]
    
    Since $a_i$ and $b_i$ are the only entries of $\ba$ and $\bb$ which disagree, all other routes of $\phi(\ba)$ and $\phi(\bb)$ are the same.
    That is,
    \[
        \phi(\ba) \setminus \phi(\bb) = \{\alpha_{1,i}\beta_{i,j}\gamma_{j,2}\}~\text{ and }~\phi(\bb) \setminus \phi(\ba) = \{\alpha_{2,i}\beta_{i,k}\gamma_{k,2}\}.
    \]
    Since $\alpha_{1,i}\beta_{i,j}\gamma_{j,2} <_i^{\cw} \alpha_{2,i}\beta_{i,k}\gamma_{k,2}$, it follows that $\phi(\ba) \lessdot \phi(\bb)$.
\end{proof}

\end{document}
